\section{Discussion}\label{sec:v8-discussion}

The v8 manuscript separates three questions that are often collapsed in
long-document measurement. Does the root score agree with the expert target?
Do the intermediate states preserve the evidence they represent? And, if the
intermediate evidence is measured through a proxy, how much uncertainty remains
after sampling true-oracle nodes? For political manifestos, that separation is
not cosmetic. It determines how to read the current results and how to spend
the next unit of expert attention.

\subsection{Reading the Certificate}

The finite-sample certificate in
Equation~\eqref{eq:v8-finite-sample-certificate} has four terms:
\[
  |G_{\mathrm{meth}}|
  \le
  C_{\mathrm{meth}}\widehat{\Delta}_{\mathrm{HT}}
  +
  B_{\mathrm{cal}}
  +
  B_{\mathrm{est}}
  +
  B_{\mathrm{clip}} .
\]
The first term is the sampled estimate of local distortion on the realized
manifesto tree. It answers a direct audit question: among the leaves,
re-summarized states, and merges that the pipeline actually used, how far did
the produced policy states move from their represented raw spans? A low root
error with a high \(\widehat{\Delta}_{\mathrm{HT}}\) is not the same evidence
as a low root error with a low \(\widehat{\Delta}_{\mathrm{HT}}\). The former
may be a successful prediction whose mechanism is not locally faithful.

\(B_{\mathrm{cal}}\) is different. It is the gap between the accessible judge
used for local evidence and the trusted policy oracle. More sampled nodes do
not remove a calibration mismatch. If quasi-sentence aggregates, teacher
traces, or LLM readouts measure a different construct from the expert-survey
dimension, then local labels can be plentiful and still point at the wrong
target. In manifesto work, this is why coded emphasis scores and
expert-position scores must be kept separate unless a calibration exercise
links them.

\(B_{\mathrm{est}}\) is the term that expert sampling can shrink. It reflects
design uncertainty around the finite population of tree units. More sampled
leaves and merges reduce this term when propensities are logged and positive.
\(B_{\mathrm{clip}}\) is the price of variance control: clipping large
inverse-propensity weights can make the estimate stable, but the induced bias
must be reported rather than hidden.

The method constant \(C_{\mathrm{meth}}\) is also substantive. A downstream
procedure that is sensitive to small score changes, or to rank reversals near a
decision threshold, should use a larger transport constant than a procedure
that only needs coarse ordering. The certificate therefore does not say that
one tree is certified for every downstream use. It certifies a tree, an oracle,
an audit design, and a method-specific tolerance.

\subsection{Root Labels and Node Labels}

Root labels and node labels do different work. Root expert labels define and
evaluate the document-level political measurement: a party's economic-policy
placement, immigration position, or decentralization stance. Without root
labels, the manuscript would have no external test of the final task. Node
labels ask a narrower question: whether a leaf, state, or merge preserved the
rubric-relevant evidence in the span it represents.

This difference matters for interpreting the current Manifesto/Benoit result.
The reported correlations are root-observed evidence. They show that the C-Tree
root state can match strong document-level benchmarks on held-out manifestos.
They do not by themselves prove C1/C2/C3 on internal nodes. The local diagnostic
figures are useful because they identify where the proxy channel is least
aligned, but they are not a completed human node-level certificate.

The next measurement layer should therefore be designed as a node survey over
the realized tree. Root labels continue to anchor the expert target. Sampled
node labels estimate local distortion and proxy bias. Quasi-sentence codings
are attractive because they are already span-grounded political evidence, but
they must be aggregated under a declared rule for the active dimension. A
decentralization audit should not reuse an economic-policy aggregation rule,
and a universal-summary audit should report whether one state is being asked
to preserve incompatible evidence geometries.

\subsection{When Local Labels Pay}

Local labels pay when three conditions hold. First, the target has meaningful
local evidence. Manifesto policy scores usually do: tax pledges, welfare
commitments, immigration restrictions, European integration claims, and
regional-autonomy statements appear in identifiable spans. Second, the tree
state has enough capacity to carry cross-span interactions, such as a tax cut
qualified by a later welfare promise or a devolution claim limited by central
budget authority. Third, the local labels measure the same construct as the
root target.

Under those conditions, node labels are more than cheaper labels. They tell the
learner what to preserve before the root score is read. They also tell the
auditor where the root score is fragile. If changing leaf size leaves the root
estimate stable, smaller leaves create more auditable sites without changing
the estimand. If a local diagnostic gap is concentrated in one dimension, as it
is for decentralization in the current artifacts, the next label should usually
go to that dimension's nodes rather than to another random root score.

Local labels pay less when the bottleneck is calibration rather than sampling.
If the accessible local target is misaligned with the expert-survey dimension,
adding more of it shrinks \(B_{\mathrm{est}}\) while leaving
\(B_{\mathrm{cal}}\) large. They also pay less when the target is intrinsically
global in a way the state map cannot localize. A rhetorical style judgment or a
party-brand interpretation may depend on document-wide framing rather than on
span evidence that can be preserved and merged. In that case the audit should
show high C1 or C3 distortion, and the certificate should widen rather than
manufacture confidence.

\subsection{Failure Modes}

The most useful failure mode is a disagreement between root accuracy and local
faithfulness. A tree can predict the right root score while losing the evidence
an expert would want to inspect. For measurement, that is not harmless: a
researcher may use the score in downstream inference, compare parties across
countries, or study which claims drove a placement. Local states make those
uses inspectable only if they remain faithful.

The second failure mode is state sharing across incompatible dimensions. The
universal-summary run asks one policy state to support six readouts. That is
efficient when dimensions share evidence, but it can dilute a dimension whose
evidence geometry is different. Decentralization is the warning case because it
depends on institutional distinctions, country-specific arrangements, and
boundary-spanning qualifications that are easy to compress away. The fix is not
necessarily a larger model. It may be a dimension-specific state map, higher
local-law weight for that dimension, or a sampling design that oversamples its
high-risk nodes.

The third failure mode is proxy overconfidence. A dense LLM judge can provide
smooth local losses everywhere, but if its errors are correlated with the
same political distinctions that experts care about, the local channel can
train the state map toward the wrong geometry. This is exactly why
Equation~\eqref{eq:v8-corrected-node-loss} uses sampled \(f^\ast\) nodes to
estimate and correct proxy bias.

\subsection{Deployment Guidance}

A manifesto deployment should declare the measurement target first. The target
is not ``summarize the manifesto''; it is a policy dimension, scale, and expert
oracle. The state rubric, root scorer, local labels, and calibration set should
all be written against that target.

The next step is to run the root-observed benchmark and the local diagnostics
together. Root correlation says whether the final score is useful. Local gaps
say where the mechanism needs audit. If the root score is strong and local
gaps are small, the certificate can focus on sampling enough nodes to tighten
\(B_{\mathrm{est}}\). If the root score is strong but local gaps are large,
the right response is targeted node labeling, not celebration of the root
number. If both root and local evidence are weak, the state map or target
rubric needs redesign.

Finally, the reported result should include the tier of evidence. A
root-observed benchmark supports a claim about held-out document scores. A
teacher-trace diagnostic supports a claim about proxy behavior. A sampled
node-level certificate supports a claim about realized local-law distortion.
Those tiers can reinforce one another, but they should not be collapsed. The
point of C-TreePO is to make the measurement stack inspectable enough that the
reader can tell which claim has actually been earned.
